\chapter*{图片致谢}

本书中的图示、地图和图表均为原创；基于已发表数据绘制的地图和图表，在图注或本列表中注明来源。照片依各自的自由许可使用，谨向其作者致谢；每张照片的完整来源链接和许可证网址，列于本书代码库中的 \texttt{images/book3/CREDITS.md}。

\bigskip
除照片和原创图示外，部分图为 AI 生成的插图，由 OpenAI 图像生成功能依据本书编辑撰写的提示词生成，并在收录前经过科学准确性审阅。这些插图只表现通用场景，从不表现真实地点、真实人物或真实标本。每张生成图像的完整提示词，记录于代码库中的 \texttt{images/book3/ai/PROMPTS.md}。

\bigskip
\noindent 照片与数据，按章列出：
\begin{itemize}\small
  \item 第~1 章（球面天文学）：帕拉纳天文台上空的星迹，R.~Wesson/ESO，CC BY 4.0。
  \item 第~2 章（望远镜、探测器与测光）：詹姆斯·韦伯空间望远镜主镜的六块镜段，低温测试前摄，NASA/MSFC/David Higginbotham，公有领域。
  \item 第~3 章（比较行星学）：按比例绘制的类地行星，NASA；金星由麦哲伦号雷达成像，NASA/JPL；火星由海盗 1 号轨道器拍摄，NASA/USGS；木卫二，NASA/JPL/DLR；木卫一，NASA/JPL/University of Arizona；均为公有领域。
  \item 第~4 章（地质图）：William Smith，\emph{A Delineation of the Strata of England and Wales, with part of Scotland}，1815；公有领域（维基共享资源）。
  \item 第~5 章（剖面图）：Hutton 在 Siccar Point 的不整合面，PawełS，CC BY-SA 4.0（维基共享资源）。
  \item 第~6 章（野外与实验室方法）：正交偏光下薄片中的黑云母花岗岩，Darla Sondrol，CC0（维基共享资源）。
  \item 第~7 章（造岩矿物）：橄榄石、石英、黑云母、铁铝榴石和角闪石标本，Robert M.~Lavinsky，CC BY-SA 3.0；拉长石，Twyla Baker，CC0；方解石解理菱面体，Giovanna Canu 与 Eva Santini，CC BY 4.0（均为维基共享资源）。
  \item 第~8 章（光学矿物学）：单偏光及正交偏光下的辉长岩薄片，Chmee2，CC BY-SA 3.0；作为学生的 Auguste Michel-Lévy，公有领域（维基共享资源）。干涉色依据 CIE 1931 色匹配函数和 CIE 标准光源 D65 计算，\copyright{} CIE，CC BY-SA 4.0。
  \item 第~9 章（分层的地球）：Andrija Mohorovičić，公有领域；1932 年的 Inge Lehmann，照片由 Even Neuhaus 拍摄，公有领域（维基共享资源）。PREM：A.~M.~Dziewonski 与 D.~L.~Anderson（1981），据 ObsPy 中的列表。
  \item 第~10 章（地球热收支）：全球热流数据来自 J.~H.~Davies 与 D.~R.~Davies（2010），\emph{Solid Earth}，CC BY 3.0；地幔绝热线来自 T.~Katsura（2022），\emph{JGR Solid Earth}，CC BY。
  \item 第~11 章（均衡与重力异常）：Vening Meinesz 摆仪装置，Broodt，CC BY-SA 4.0（维基共享资源）。高程数据来自 ETOPO1（NOAA 国家环境信息中心），按一度间隔取样。
  \item 第~12 章（球面上的板块运动学）：一座永久性卫星定位站，Alpsdake，CC BY-SA 4.0（维基共享资源）。板块边界来自 PB2002（P.~Bird，2003；GeoJSON 由 H.~Ahlenius 制作，ODC-By 1.0）；海岸线来自 Natural Earth（公有领域）。
  \item 第~13 章（裂谷、洋脊与被动大陆边缘）：冰岛辛格维利尔国家公园的张开裂隙，Olga Ernst，CC BY-SA 4.0；加拉帕戈斯裂谷的枕状熔岩，NOAA Okeanos Explorer Program，CC BY 2.0（均为维基共享资源）。深度来自 ETOPO1（NOAA 国家环境信息中心），洋底年龄来自 M.~Seton 等人（2020），EarthByte，CC BY 4.0。
  \item 第~14 章（活动大陆边缘）：从国际空间站所见喷发中的 Sarychev 峰，NASA，公有领域（维基共享资源）。地震来自 USGS ANSS 综合目录（公有领域）；火山来自史密森学会全球火山活动计划；板块边界来自 PB2002（P.~Bird，2003；GeoJSON 由 H.~Ahlenius 制作，ODC-By 1.0）；海岸来自 Natural Earth（公有领域）；洋底年龄来自 M.~Seton 等人（2020），EarthByte，CC BY 4.0。
  \item 第~15 章（热点与板内火山作用）：基拉韦厄 Pu\textquoteleft u \textquoteleft O\textquoteleft o 的熔岩喷泉，G.~E.~Ulrich，美国地质调查局，公有领域；马特兰的西高止山脉，Nicholas（Nichalp），CC BY-SA 2.5（均为维基共享资源）。热点位置、轨迹方位角与速率来自 W.~J.~Morgan 与 J.~P.~Morgan（2007）电子补充材料中的表~1；深度来自 ETOPO1（NOAA 国家环境信息中心）；板块边界来自 PB2002（P.~Bird，2003；GeoJSON 由 H.~Ahlenius 制作，ODC-By 1.0）；海岸来自 Natural Earth（公有领域）。
  \item 第~16 章（应力、应变与流变学）：用于测定岩石强度的高压三轴仪，Darryl Peroni，CSIRO，CC BY 3.0（维基共享资源）。湿石英岩的位错蠕变参数（Gleason 与 Tullis，1995）和干橄榄岩的位错蠕变参数（Hirth 与 Kohlstedt，2003），由 M.~A.~Lopez-Sanchez 在 \emph{strength\_envelopes} 中汇编（MPL-2.0）。
  \item 第~17 章（断层、褶皱与面理）：旧金山 Corona Heights 的擦痕面，Dextersinister，CC BY-SA 4.0；内华达 Keystone 逆冲断层的断层角砾岩，Jstuby，公有领域；康沃尔 Millook 的人字形褶皱，Smalljim，CC BY 3.0；坎塔布里亚 Potes 附近带轴面劈理的背斜，Mikenorton，CC BY-SA 3.0（均为维基共享资源）。
  \item 第~18 章（地震震源）：1906 年马林县破裂错开的土地与围栏，\emph{Popular Science Monthly} 第 69 卷（1906），公有领域；圣安德烈亚斯断层错开的 Wallace 溪谷，R.~E.~Wallace，美国地质调查局，公有领域（均为维基共享资源）。地震来自 USGS ANSS 综合目录（公有领域）；2011 年东北地震的震源机制来自 USGS W 相矩张量解（公有领域）。
  \item 第~19 章（风化作用与黏土）：比利牛斯山脉 Targassonne 的风化花岗岩与花岗砂，Pinpin，CC~BY~2.5；澳大利亚 Weipa 高岭质砂岩上的铝土矿，Werner Schellmann，CC~BY-SA~2.5；扫描电子显微镜下的高岭石晶体，美国地质调查局，公有领域（均为维基共享资源）。离子半径来自 R.~D.~Shannon（1976）。
  \item 第~20 章（侵蚀与搬运）：2018 年 3 月 4 日，MODIS 仪器在 Terra 卫星上拍摄的密西西比河洪水羽流，NASA Earth Observatory 图像，Joshua Stevens 制作，公有领域（维基共享资源）。
  \item 第~21 章（碎屑沉积物及其沉积环境）：内华达 Aztec 砂岩中的交错层理，Chip Carroon，土地管理局，CC~BY-SA~2.0；Landsat~7 拍摄的勒拿河三角洲，美国地质调查局，公有领域；具有 Bouma 序列各段的浊积岩，Becke-Oese，Jo Weber，CC~BY-SA~3.0（均为维基共享资源）。
  \item 第~22 章（碳酸盐、硅质与蒸发岩沉积）：从国际空间站所见的波拉波拉岛，NASA 约翰逊航天中心，公有领域；\emph{Emiliania huxleyi} 的一个细胞，Alison R.~Taylor，北卡罗来纳大学威尔明顿分校显微镜设施，CC~BY~2.5；\emph{Orbulina universa}，H.~Takagi、K.~Kimoto、T.~Fujiki、H.~Saito、C.~Schmidt、M.~Kucera 与 K.~Moriya，CC~BY~4.0；麦克默多湾的硅藻，Gordon T.~Taylor，石溪大学，NOAA Corps Collection，公有领域（均为维基共享资源）。
  \item 第~23 章（成岩作用与沉积盆地）：Potsdam 砂岩中的石英加大边，M.~C.~Rygel，CC~BY-SA~4.0（维基共享资源）。海平面叠合曲线来自 R.~M.~Spratt 与 L.~E.~Lisiecki（2016），NOAA NCEI 古气候世界数据服务中心；压实参数来自 J.~G.~Sclater 与 P.~A.~F.~Christie（1980），据 pyBacktrack（EarthByte）中的列表。
  \item 第~24 章（地层年代表）：Zumaia 塞兰阶 GSSP 的金钉，Josugoni，CC~BY-SA~4.0（维基共享资源）；Trinidad Lake 州立公园的白垩纪--古近纪界线，Nationalparks，CC~BY-SA~2.5（维基共享资源）。国际年代地层表的底界年龄与颜色，国际地层委员会，CC~BY~4.0。
  \item 第~25 章（放射性测年方法）：阴极发光下的锆石颗粒，Emmanuel Roquette，CC~BY-SA~4.0；热电离质谱仪，Radiogenic，CC0（均为维基共享资源）。IntCal20 校准曲线来自 P.~J.~Reimer 等人（2020），\emph{Radiocarbon} 第 62 卷，CC~BY~4.0，由 M.~Blaauw 随 rintcal 软件包发布。
  \item 第~26 章（大气与海洋的结构）：在白令海峡回收 CTD 与采水器，NASA 戈达德太空飞行中心，公有领域（维基共享资源）。世界海洋图集 2023 的温度与盐度（NOAA 国家环境信息中心）。
  \item 第~27 章（辐射与热量输送）：深空气候观测站拍摄的被阳光照亮的地球，NASA 戈达德太空飞行中心，公有领域（维基共享资源）。NCEP-NCAR 再分析资料第 1 版的辐射、地表通量与降水（E.~Kalnay 等人，1996），由美国科罗拉多州博尔德 NOAA PSL 提供，网址 https://psl.noaa.gov。
  \item 第~28 章（大气环流）：阿拉斯加湾中的一个温带气旋，NASA MODIS/LANCE，公有领域；GOES-11 云图所示的热带辐合带，NASA Earth Observatory，公有领域（均为维基共享资源）。NCEP-NCAR 再分析资料第 1 版的风场与海平面气压（E.~Kalnay 等人，1996），由美国科罗拉多州博尔德 NOAA PSL 提供，网址 https://psl.noaa.gov。
  \item 第~29 章（海洋环流）：由海表温度所见的墨西哥湾流，NOAA NESDIS，公有领域；正在投放的深海剖面浮标，未署名 NOAA 摄影师，公有领域（均为维基共享资源）。
\end{itemize}
